\section{Dynamic Speculative Planning with User Controllability}
Dynamic Speculative Planning begins with a simple intuition: once $\approxi$ has produced a few actions, it is often possible to forecast how many additional actions it will generate correctly before $\target$ must intervene. Rather than selecting $k$ heuristically or relying on an expensive offline model, we learn to predict it \emph{online} while the system is running. After one speculative execution breaks either due to a detected mismatch or when a previously the step is exhausted, the system logs a training pair (partial trajectory, $k$) for training while at the same time utilize the latest checkpoint of the predictor to predict the next $k$. These training pairs are collected to train a small DistilBERT-based regressor whose training cost can be negligible compared with inference cost of a single LLM call. To avoid the training time impacting the speed of the system, the predictor training is done asynchronously: for every newly collected pairs of datapoints, an independent thread fine‑tunes the predictor to convergence and then writes the updated checkpoint to disk. The orchestration layer always utilizes the most recent checkpoint, so learning progresses without ever blocking agent planning. 

The online, multi-threaded training regimen offers two key advantages: (1) \textit{Real‑time adaptation}: the predictor is continuously updated with the most recent user interactions, keeping it aligned with the current usage distribution. (2) \textit{Prompt cost savings}: the system begins reducing expenses as soon as new data arrive, without waiting to accumulate and curate a large offline dataset.


\subsection{Speculative‑Step Control as Value Predictions}
\label{sec:spec_step_control}

As mentioned above, achieving dynamic speculative planning requires predicting the $k$ for each (partial) trajectory.
Moreover, this predictor must be updated online to enable continuous adaptation.
These requirements make online reinforcement learning a suitable framework.
Specifically, we can formalize speculative step prediction as a state-value prediction problem in reinforcement learning and apply Temporal-Difference (TD) learning~\citep{sutton2018reinforcement} to update a state-value function in an online and incremental manner.

Formally, let $\Dist(X)$ be the space of all probability distributions supported over the set $X$.
Consider a Markov decision process (MDP), $M=(\States, \Actions, \Transition, p_0, R, \gamma)$, where the state space $\States$ is a set of token sequences, the action space $\Actions$ is a set of speculative steps, $\Transition: \States \times \Actions \rightarrow \Dist(\States)$ is the transition function, the initial state distribution $p_0$ is a probability distribution over the prompt dataset $Q$ (\emph{i.e.}, $p_0 \in \Dist(Q)$), $R: \States \times \Actions \rightarrow \mathbb{R}$ is the reward function, and $\gamma \in [0,1]$ is the discount factor.

An approximation agent $\approxi$ interacts with the MDP based on a policy $\pi: \States \rightarrow \Dist(\Actions)$. Specifically, for each episode, the agent starts from a prompt $s_0 \sim p_0(\cdot)$.
At each time-step $t$, it observes the state $s_t \in \States$, takes an action (\emph{i.e.}, a speculative step) $a_t \sim \pi(\cdot|s_t)$, transitions to a new state $s_{t+1}$ and receives a scalar reward $r_{t+1} = R(s_t, a_t)$.
The episode ends when an approximation-target mismatch is met.
In our case, $s_{t+1}$ is simply a concatenation of the token sequences of $s_t$ and $a_t$, \emph{i.e.}, $s_{t+1} = s_t | a_t$. We also set $\gamma=1$.
After an episode ends, we get a complete trajectory and it is formally defined as $\tau = (s_0, a_0, r_1, s_1, \dots, s_T)$, where $T$ is the final time-step. 
Define return $G_t$ over $\tau$ as the total (discounted) reward from time-step $t$: $G_t = \sum_{i=t}^{T-1} \gamma^{i-t} R(s_i, a_i).$
A state-value function is defined as the expected return under policy $\pi$, $V_{\pi}(s) = \E_{\pi}[G_t | s_t=s].$

The reward function is designed so that the return $G_t$ corresponds to the speculative step at time-step $t$.
Specifically, if the speculative step $a_t$ (generated by $\approxi$) differs from the one generated by $\target$, we set $R(s_t, a_t) = 1$; otherwise, $\approxi$ gets zero reward and the episode ends immediately.
Under this reward setting, the state-value $V_{\pi}(s)$ is essentially the expected speculative step of state $s$.

In practice, $V_{\pi}$ is usually approximated by a neural network $V_{\theta}$, where $\theta$ denotes the network weights.
To learn a good approximation of $V_{\pi}$, we apply TD learning and minimize the following loss function:
\begin{align}\label{eq:td_loss}
L_{\theta} = \E_{\tau \sim \pi} [(G_t^\lambda - V_{\theta}(s_t))^2],
\end{align}
where $G_t^\lambda$ is known as $\lambda$-return~\citep{sutton2018reinforcement}.
Formally, it is defined as
\begin{align}
G_t^\lambda = (1-\lambda) \sum_{n=1}^{T-t-1} \lambda^{n-1} G_{t: t+n} + \lambda^{T-t-1} G_t,
\end{align}
where $G_{t: t+n} = r_{t+1} + \gamma r_{t+2} + \cdots + \gamma^{n-1} r_{t+n}+\gamma^n V_{\theta}(s_{t+n})$ and $\lambda \in [0,1]$.
The usage of $\lambda$-return offers a flexible way to achieve a better bias-variance tradeoff by interpolating between Monte Carlo return $G_t$ (when $\lambda=1$) and 1-step TD target $r_{t+1} + \gamma V_{\theta}(s_{t+1})$ (when $\lambda=0$).
It is worth noting that when $\lambda=1$ and $\gamma=1$, TD learning in this context reduces to online supervised learning since $G_t^1 = G_t = k$, thus our learning framework is a more generalized framework.
During training, since rewards depend on the speculative steps generated by $\target$ and $\target$ runs slower than $\approxi$, we wait until $\target$ finishes and then compute rewards.
However, during deployment, a trained state-value function $V_{\theta}$ can be used to predict the speculative step immediately, without waiting for $\target$ to finish.


\subsection{Multi-Thread Online Learning}

To minimize inference latency and update overhead, our system adopts a fully asynchronous learning framework that decouples predictor inference and training from the core execution pipeline: (1) \textbf{asynchronous inference for parallel execution}: To mitigate the impact of predictor latency on system throughput, we perform speculative step prediction asynchronously at the start of each episode. The predictor runs concurrently with $\approxi$'s planning step, and typically completes before $\approxi$'s current step finishes. This parallelism effectively hides prediction latency, ensuring uninterrupted speculative planning.
%and preserving end-to-end efficiency.
% We predict the forthcoming speculative step $k$ at the beginning of each speculative episode. Although the predictor model is substantially more efficient than invoking $\approxi$, its inference latency could still impact overall system performance if implemented sequentially. To address this potential bottleneck, we employ an asynchronous inference strategy that effectively masks the prediction latency. Specifically, we invoke the predictor asynchronously in parallel with the thread generating $\approxi$'s next step. Since predictor inference typically completes more rapidly than a single agent step generation, the predicted speculative step $k$ is generally available by the time the agent completes its current step. This parallel execution architecture ensures that speculative planning proceeds without interruption, eliminating additional delays that would otherwise accumulate during sequential prediction.
(2) \textbf{asynchronous training for continuous adaptation}: To adapt to evolving task distributions, the predictor is updated online via a decoupled training-inference architecture. We maintain two models: a \emph{training} predictor that undergoes continuous updates, and an \emph{inference} predictor that serves all runtime inference requests. 
Upon completion of each speculative episode, an asynchronous thread is launched to train the predictor on the collected training dataset, which operates independently of the core planning pipeline. Once the training finishes, updated weights are immediately synchronized from the training predictor to the in-memory inference predictor.
% In addition to parallel inference, the predictor is continuously updated online to adapt to the evolving task distribution. 
% We decouple the training and inference of the predictor by maintaining two separate models: a \emph{training} predictor, which undergoes continuous updates, and a \emph{inference} predictor, which serves all inference requests during execution.

 % This approach eliminates the overhead of disk I/O and ensures that inference always benefits from the most recent updates without interrupting the core execution.

% \paragraph{Robustness and Seamless Recovery} To further enhance robustness, our system is designed to support seamless recovery in the event of network disruptions, task interruptions, or user-initiated restarts. After the completion of each training thread and synchronization, the latest model parameters are also persisted to disk. In addition, at the end of each task, the corresponding training data is also stored. This enables a smooth restarted training from last checkpoint, eliminating the need to retrain from scratch.


% This multi-threaded design offers two principal benefits. First, it supports real-time adaptation: the predictor is continually refined with fresh data, allowing it to remain aligned with evolving user and task dynamics. Second, it is designed to largely avoid introducing latency into the execution pipeline: by decoupling inference and training from agent execution, the system maximizes responsiveness while incrementally improving prediction quality. Finally, it enables robust and seamless recovery: persistent storage of predictor states and training episodes allows the system to resume online learning without retraining, even after interruptions or restarts.%